%% ma: L11-B14-0010
%% lop: 11
%% chuong: 4
%% bai: 14
%% dang: 11.14.2
%% loai: TN4
%% muc: TH
%% dap_an: "C"
%% nguon: "tai-lieu-hoc-tap-mon-toan-lop-11-theo-chuong-trinh-giao-duc-pho-thong-2018.pdf (Lớp Toán Thầy Duy, Phạm Lê Duy, THPT Chu Văn An), Chương 4 Bài 5, Phần 1, Câu 5 (tr.526)"
%% kien_thuc: "Lớp 11 Bài 14 (phép chiếu song song)"
%% trang_thai: cho-duyet
%% ly_do: "Dạng mới đề xuất (11.14.x), chờ giáo viên duyệt dạng; tài liệu gốc không có lời giải mẫu, đáp án do kho tự giải và kiểm bằng sympy."
\begin{ex}
Qua phép chiếu song song lên mặt phẳng $(P)$, hai đường thẳng $a$ và $b$ có hình chiếu là hai đường thẳng song song $a^{\prime}$ và $b^{\prime}$. Khi đó:
\choice{$a$ và $b$ phải song song với nhau.}{$a$ và $b$ phải cắt nhau.}{\True $a$ và $b$ có thể chéo nhau hoặc song song với nhau.}{$a$ và $b$ không thể song song.}
\loigiai{\sloppy Nếu $a,b$ cắt nhau thì hình chiếu cắt nhau hoặc trùng nhau, không thể song song nên loại. Vậy $a,b$ song song hoặc chéo nhau (cả hai đều có thể xảy ra), đáp án C.}
\end{ex}
